\subsection*{Defining Shapes with \texttt{mdspan}}

The keyword \co@mdspan@ stands for multidimensional span. It describes the extents of a shape without allocating its elements. For example, a shape \co@[7, 8]@ has rank 2 and describes 56 positions; choosing \co@f32@ or \co@s32@ for those positions is a separate decision.

You can explicitly define \texttt{mdspan} variables as follows:

\begin{lstlisting}[language=croqtile,caption={Declaring a shape}]
mdspan s0 : [7, 8]; // Defines a 2D shape with dimensions [7, 8]
mdspan<1> s1 : [3]; // Defines a 1D shape with dimensions [3]
\end{lstlisting}

The colon introduces the shape initializer, and square brackets enclose its dimensions. Thus \co@s0@ has seven rows and eight columns, while \co@s1@ has one dimension of size 3. A named shape can be reused in later buffer declarations.

You can specify the \textbf{rank} by placing \texttt{\textless{}rank\textgreater{}} after \texttt{mdspan}. The compiler checks that this rank matches the initialization expression and reports a compilation error if they differ:

\begin{lstlisting}[language=croqtile,caption={A rank mismatch}]
mdspan<3> s2 : [64, 32]; // error: the rank of mdspan is inconsistent
\end{lstlisting}

Because the initializer identifies a shape, the compiler can infer the \co@mdspan@ type when the keyword is omitted:

\begin{lstlisting}[language=croqtile,caption={Omitting the mdspan keyword}]
s3 : [7, 8, 9]
\end{lstlisting}

Here \co@s3@ has rank 3 and dimensions 7, 8, and 9. Every named shape requires an initializer, so the inferred form is often sufficient.
